\begin{textsample}{POS Dim 1 – ap – Score 42.00 – ap/ap\_ef\_3.txt}  \label{ex:f1_pos_174}
4.2.2 Desenvolvimento e Aprendizagem Ao longo do Ensino Fundamental Anos Iniciais e Finais, a progressão do conhecimento ocorre pela consolidação das aprendizagens anteriores e pela ampliação das práticas de linguagem e da experiência estética e intercultural das crianças, considerando tanto seus interesses e suas expectativas quanto o \textbf{que} ainda precisam aprender.
Ampliam -se a autonomia intelectual, a compreensão de normas e os interesses pela vida social, o \textbf{que} lhes possibilita lidar com sistemas mais amplos, \textbf{que} \textbf{dizem} respeito às relações dos sujeitos entre si, com a natureza, com a história, com a cultura, com as tecnologias e com o ambiente.
Além desses aspectos relativos à aprendizagem e ao desenvolvimento, na elaboração dos currículos e das propostas pedagógicas devem ainda ser consideradas medidas para assegurar aos alunos \textbf{um} percurso contínuo de aprendizagens entre as duas fases do Ensino Fundamental, de modo a promover \textbf{uma} maior integração entre elas.
Afinal, essa transição se caracteriza por mudanças pedagógicas na estrutura educacional, decorrentes principalmente da diferenciação dos componentes curriculares.
Como bem destaca o Parecer CNE/CEB nº 11/2010, “ os alunos, ao mudarem do professor generalista dos anos iniciais para os professores especialistas dos diferentes componentes curriculares, costumam se ressentir diante das muitas exigências \textbf{que} têm de atender, feitas pelo grande número de docentes dos anos finais ” ( BRASIL, 2010 ).
Realizar as necessárias adaptações e articulações, tanto no 5º quanto no 6º ano, para apoiar os alunos nesse processo de transição, pode evitar ruptura no processo de aprendizagem, garantindo -lhes maiores condições de sucesso.
Os estudantes dessa fase inserem -se em \textbf{uma} faixa etária \textbf{que} corresponde à transição entre infância e adolescência, marcada por intensas mudanças decorrentes de transformações biológicas, psicológicas, sociais e emocionais.
Nesse período de vida, como bem aponta o Parecer CNE/CEB nº 11/2010, ampliam -se os vínculos sociais e os laços afetivos, as possibilidades intelectuais e a capacidade de \textbf{raciocínios} mais abstratos.
Os estudantes tornam -se mais capazes de \textbf{ver} e avaliar os fatos pelo ponto de vista do outro, exercendo as capacidades de empatia e resiliência, “ importante na construção da autonomia e na aquisição de valores morais e éticos ” ( BRASIL, 2010 ).
As mudanças próprias dessa fase da vida \textbf{implicam} a compreensão do adolescente como sujeito em desenvolvimento, com \textbf{singularidades} e formações \textbf{identitárias} e culturais próprias, \textbf{que} \textbf{demandam} práticas escolares diferenciadas, capazes de contemplar suas necessidades e diferentes modos de inserção social.
Conforme reconhecem as DCN, é frequente nessa etapa.
A aprendizagem é \textbf{um} processo contínuo \textbf{que} ocorre durante toda a vida do indivíduo, desde a mais tenra infância até a mais avançada velhice.
Normalmente \textbf{uma} criança deve aprender a andar e a falar ; depois a ler e escrever, aprendizagens básicas para atingir a cidadania e a participação ativa na sociedade.
Já os adultos precisam aprender habilidades ligadas a algum tipo de trabalho \textbf{que} lhes forneça a satisfação das suas necessidades básicas, algo \textbf{que} lhes não somente garanta o seu sustento, mas seu desenvolvimento como cidadão. \textbf{Contudo}, a educação não acontece em \textbf{um} prédio, mas em todo e qualquer lugar, \textbf{uma} vez \textbf{que} as pessoas aprendem na interação, conforme esclarece Vygotsky ( 1998 ), por meio de mediações com alguém mais experiente como \textbf{um} adulto e também na relação com os objetos.
O processo educativo por fazer parte da vida de todos os seres, pois “ a educação aparece sempre \textbf{que} surgem formas sociais de condução e controle da aventura de ensinar e aprender ” ( BRANDÃO, 1998, p. 26 ).
Corrêa ( 2006 ) afirma \textbf{que} o educador precisa \textbf{atentar} -se para o futuro e, portanto, às novas necessidades socioculturais, o \textbf{que} \textbf{exige} de o professor buscar novas formas de entender e pensar o seu papel nesse contexto histórico.
Nogueira ( 2007 ) explica \textbf{que} a formação continuada, como necessidade e obrigação, \textbf{permeia} a evolução profissional de educadores, e torna -se cada vez mais essencial na teia complexa da vida social e acadêmica de docentes.
O \textbf{que} vem a ser aprendizagem?
A aprendizagem é fruto da história de cada \textbf{sujeito} e das relações \textbf{que} ele \textbf{consegue} estabelecer com o conhecimento ao longo da vida.
No entanto, esse processo não \textbf{depende} somente do indivíduo, é \textbf{um} processo \textbf{que} envolve as interações \textbf{que} a pessoa realiza com o outro, com a família, amigos etc., ou em situações formais, como com o professor, seus colegas, livros dentre outros.
Segundo Fernandez ( 2001 ) a família também é responsável pela aprendizagem da criança, já \textbf{que} os pais são os primeiros a ensinar e os mesmos determinam algumas modalidades de aprendizagem aos filhos.
Para Vygotsky ( 1989 ) a aprendizagem da criança começa muito antes da aprendizagem escolar \textbf{que} nunca parte do zero.
A aprendizagem da criança na escola ocorre com o estabelecimento de relações \textbf{que} se conectam com os conhecimentos anteriores da criança.
Fernandez ( 1996 ) aponta \textbf{que} a aprendizagem é colocada como apropriação, reconstrução do conhecimento do outro a partir do saber pessoal.
Resulta da interação entre as estruturas do pensamento e o meio cultural \textbf{que} \textbf{necessita} ser compreendido.
Davis e Oliveira ( 1993, p. 20-21 ) destacam a importância da interação e socialização da criança em sua aprendizagem, apontam \textbf{que} : A aprendizagem é o processo através do qual a criança se apropria ativamente do conteúdo da experiência humana, daquilo \textbf{que} o seu grupo social conhece.
Para \textbf{que} a criança aprenda, ela \textbf{necessitará} interagir com outros seres humanos, especialmente com os adultos e com outras crianças mais experientes.
Acompanhar a evolução social e o contexto em \textbf{que} os alunos se encontram inseridos.
Há ainda a necessidade de destacar \textbf{que} as propriedades da aprendizagem são paralelas às da evolução tendo em vista \textbf{que} toda estrutura educacional está organizada com o intuito de, primeiramente, promover a aprendizagem e o desenvolvimento do ser humano.
Dessa forma, a busca pela aprendizagem deve ser estabelecida em alguns tipos de contingências, sendo assim é necessário observar o saber como, relacionado ao conceito operacional, o saber sobre ( análise das instituições de ensino e o contexto \textbf{atual} ).
O desenvolvimento e a aprendizagem requerem mudanças, relativamente permanente no comportamento, resultante do conhecimento e da experiência.
Moreira ( 2003 ) explica \textbf{que} a aprendizagem significativa é aquela \textbf{que} ocorre com significado, quando está inserida em conceitos, ideias, proposições, modelos e fórmulas, \textbf{que} estimulam algo novo ao \textbf{aprendiz} e diante disso, aponta \textbf{que} o ensino deve buscar sempre a facilitação da aprendizagem, utilizando o princípio da interação e do questionamento para \textbf{que} o aluno aprenda de maneira significativa.
O professor este deve relacionar o conteúdo de maneira não arbitrária e não literal, ou seja, ele deve auxiliar os alunos a tecerem relações entre os conhecimentos aprendidos anteriormente com os novos conhecimentos.
As aprendizagens repousam dessa forma sobre \textbf{um} tripé : quem aprende, o \textbf{que} se aprende e o outro, ou seja, repousa sobre o \textbf{sujeito}, o objeto, o social.
Outro ângulo de \textbf{teoria} de aprendizagem é o da esfera \textbf{que} comporta a problemática dos significados, dos valores, do sentido da vida, estes valores são significados e constituídos pelo âmbito cultural dos grupos.
Podemos destaca \textbf{que} a escola deve propiciar educação de base com qualidade, \textbf{que} alargue e potencialize as capacidades dos alunos, para isso é necessário desenvolver o uso dos processos internos, biológicos e sociais.
A compreensão e a criatividade devem ser buscadas pela escola e pelo educador, visando assim estimular o conhecimento transferido.
É \textbf{imprescindível} \textbf{que} o aluno possa agir com autonomia, ou seja, sujeitos próprios \textbf{corresponsáveis} pelas suas aprendizagens e susceptíveis às transformações.
A aprendizagem se difunde pela vida do \textbf{homem} em todos os seus aspectos e, diante disso, aponta \textbf{que} as \textbf{teorias} de aprendizagem, buscam reconhecer a dinâmica envolvida nos atos de ensino aprendizagem, partindo do conhecimento da evolução cognitiva do \textbf{homem}, \textbf{visto} \textbf{que} o conceito de aprendizagem tem vários significados não compartilhados sendo expressos em três enfoques \textbf{teóricos} principais \textbf{que} são : o enfoque behaviorismo, cognitivista, construtivista e sócio histórica.
O enfoque behaviorista é a base \textbf{teórica} \textbf{que} fundamenta o ensino chamado tradicional, com professor atuando de forma \textbf{central} no processo de ensino, o construtivismo \textbf{que} é o suporte \textbf{teórico} para a escola nova e, sujeito é o centro do processo de aprendizado, o sócio histórico \textbf{que} aborda o processo de forma \textbf{dialética} em \textbf{que} o professor, como mediador do processo de ensino interage com o aluno de modo a potencializar seu aprendizado e desenvolvimento.
Precisamos está como olhar atento com relação ao comportamento e \textbf{que} a aprendizagem se firma na observação do desenvolvimento cognitivo, tendo como base os demais alunos.
O cognitivismo enfatiza aquilo \textbf{que} é ignorado pelo ponto de vista behaviorista, \textbf{que} seria a \textbf{cognição}, o ato de conhecer, ou seja, como \textbf{um} ser humano conhece o mundo, investigam a \textbf{percepção} e a compreensão, dentre outros processos mentais dos seres humanos, de forma científica.
Mas não podemos deixar de \textbf{lado} o papel do aluno, é \textbf{preciso} ter consciência \textbf{que} \textbf{uma} boa escola, \textbf{uma} boa metodologia é construída no dia a dia do fazer escolar e para tanto precisa de \textbf{disciplina}, atenção e \textbf{um} pensamento lógico com apoio das tecnologias disponíveis para o processo ensino-aprendizagem, levando a \textbf{um} desenvolvimento satisfatório.
O aluno não é mais \textbf{simples} receptáculo de conhecimento, ele é parte fundamental do meio escolar onde está inserido, e por tanto deve interagir com professores, coordenadores e com seus colegas de sala.
Nos dias \textbf{atuais} os \textbf{métodos} devem ser : dinâmicos, inovadores, variados e \textbf{que} leve o aluno a aprender fazendo, e \textbf{que} esteja de acordo com sua realidade.
Pois isso proporciona a quem aprende diferentes formas de \textbf{ver} e refletir sobre sua importância no mundo \textbf{atual}, através da investigação dos fatos \textbf{que} pertencem ao seu meio.
Dessa forma, pode -se colocar \textbf{que} a aprendizagem humana é determinada pela interação entre o indivíduo e o meio em \textbf{que} participam os aspectos biológicos, psicológicos e sociais.
Nesse contexto é pertinente concluir \textbf{que} é fundamental \textbf{que} a criança seja \textbf{instigada} em sua criatividade e \textbf{que} suas curiosidades sejam satisfeitas por meio de descobertas concretas, desenvolvendo a sua autoestima, criando em si \textbf{uma} maior segurança e confiança, tão necessárias à vida adulta.
É \textbf{preciso} \textbf{que} os pais se envolvam nos processos educativos dos filhos, no sentido de motivá -los afetivamente ao aprendizado.
O aprendizado formal ou a educação escolar, para ser bem-sucedida, não \textbf{depende} apenas de \textbf{uma} boa escola ou de bons programas, mas também de como a criança é tratada em casa e dos \textbf{estímulos} \textbf{que} recebe para aprender.
É \textbf{preciso} entender \textbf{que} o aprender é \textbf{um} processo contínuo e não cessa quando a criança está em casa, é \textbf{preciso} estratégias, inovação e criatividade para \textbf{que} o processo ensino-aprendizagem aconteça.
Destaca -se \textbf{que} os procedimentos da aprendizagem, deverão estar relacionados com a realidade dos seus educandos, seu contexto-social, \textbf{que} aborda o processo de forma \textbf{dialética} em \textbf{que} o professor, como mediador de processo de ensino interage com o aluno de modo a potencializar seu aprendizado e desenvolvimento.
Tem \textbf{uma} preocupação com os aspectos observáveis do comportamento.
Davis e Oliveira também seguidores da \textbf{abordagem} sócio histórica, podemos pensar na importância da motivação no contexto de aprendizagem e, \textbf{que} a motivação para aprender nada mais é do \textbf{que} o reconhecimento pelo indivíduo da necessidade de aprendizagem, sendo \textbf{que} dessa forma o professor tem como papel estabelecer de forma mediadora a interação do indivíduo com a informação, planejando atividades significativas \textbf{que} sejam motivadoras e \textbf{interessantes} de modo \textbf{que} os alunos busquem o conhecimento e a aprendizagem.
Quanto a Escola, \textbf{que} é \textbf{um} lugar de produção de saber por \textbf{excelência}.
Planejar atividades significativas em sala de aula para \textbf{que} o aluno estabeleça relações entre os conceitos cotidianos e os conceitos científicos, de modo a elaborar seus conhecimentos é \textbf{uma} das exigências para o desenvolvimento de \textbf{um} processo educativo \textbf{ancorado} em \textbf{uma} \textbf{abordagem} sócio histórica.
Pois o conhecimento é condição necessária para o domínio do \textbf{homem} sobre a natureza e sobre os outros \textbf{homens}, porém, a escola não é a única instituição educativa da sociedade, esta é caracterizada como tendo função primordial a transmissão do saber e formação dos indivíduos.
Assim, para \textbf{que} a função educativa da escola aconteça, é necessário \textbf{que} o professor e o corpo coletivo da escola estejam capacitados para entender a ampla realidade dos \textbf{sujeitos} e assim desenvolver, aprofundar e ampliar o conhecimento.
Para tanto, o professor tem \textbf{que} ter autoria, autonomia e meios de aperfeiçoar a sua capacidade de analisar e interpretar a realidade ( FRIGOTTO, 1995 ). \textbf{Um} dos principais papéis da escola é o de oferecer situações de aprendizagens, para \textbf{que} o aluno possa conhecer e fazer uso do conhecimento aprendido. \textbf{Uma} boa condição de aprendizagem requer \textbf{que} o educador considere o grupo-classe e suas possibilidades e também avaliar as condições individuais dos mesmos.
Para \textbf{que} \textbf{uma} atividade tenha bons resultados é necessário \textbf{que} o professor planeje a organização \textbf{que} utilizará na atividade \textbf{que} irá propor.
O professor deve encorajar seus alunos a pensar por si mesmos, sem querer obter deles, respostas e soluções corretas, mas \textbf{que} valorizem também as tentativas, e desafios vivenciados.
Com relação aos espaços físicos e mobiliários, \textbf{que} estes também são fatores importantes.
A escola, as salas de aulas e as aulas ministradas requerem capacidade de oferecer ensino de qualidade para tanto, faz -se necessário a utilização de espaço físico, materiais e recursos didáticos favoráveis ao processo educativo dos alunos.
A educação em sala de aula é \textbf{um} processo discursivo \textbf{sócio-histórico}, no qual os resultados, do ponto de vista da aprendizagem são determinados conjuntamente pelos esforços de professores e alunos.
Porém, o aluno é o responsável final pela sua aprendizagem ao atribuir significado aos conteúdos, mas é o professor \textbf{que}, com sua intervenção, oferece atividades \textbf{que} podem possibilitar a construção do conhecimento ( COLL ; EDWARDS, 1998 ).
A aprendizagem dos alunos se apresenta como objetivo do trabalho docente, e, por \textbf{conseguinte}, o importante é a organização apropriada das atividades escolares para \textbf{que} alcancem esse objetivo, e para dirigir tais situações de aprendizagem, é indispensável \textbf{que} o professor domine os saberes, \textbf{que} esteja mais de \textbf{uma} lição a frente dos alunos e \textbf{que} seja capaz de encontrar o essencial sob múltiplas aparências, em contextos variados.
Com relação ao processo de troca e interação necessária entre professor e aluno, o professor deve trabalhar com os alunos os significados dos conceitos científicos e o aluno, por sua vez, deve devolver ao professor o \textbf{que} ele apreendeu.
Dessa forma, professores e alunos têm responsabilidades distintas, o professor na transmissão de conteúdos e conhecimentos, na avaliação do \textbf{que} foi compreendido pelo aluno e o aluno se responsabiliza pelo estudo, participação e expressão do \textbf{que} aprendeu.
Dessa forma, o processo de ensino e aprendizado requer responsabilidades distintas e interativas.

% matched lemmas: abordagem, ancorar, aprendiz, atentar, atual, central, cognição, conseguinte, conseguir, contudo, corresponsável, demandar, depender, dialético, disciplina, dizer, estímulo, excelência, exigir, homem, identitária, implicar, imprescindível, instigar, interessante, lado, método, necessitar, percepção, permear, preciso, que, raciocínio, simples, singularidade, sujeito, sócio-histórico, teoria, teórico, um, ver
\end{textsample}
